\documentclass{article}

\begin{document}
    \begin{center}
        \Large \textbf{Challenge 5: Linux für Fortgeschrittene und Experten} \\ [6pt]
        \normalsize
        Autor: Maximilian Jakob Maag B.Sc. \\
        Version: 1.0
    \end{center}

    \section{Intro}
    Diese Challenge richtet sich an fortgeschrittene Nutzer.
    Hier geht es um tiefgehende Systemanpassungen, Sicherheit und moderne Technologien. \\\\
    Sie erhalten einen Einblick in komplexere Themenbereiche rund um Linux. \\\\

    \section{Lernziele}
    Sie lernen fortgeschrittene Linux-Konzepte kennen.
    Verschaffen Sie sich einen Überblick über moderne Technologien und Sicherheitsmechanismen.

    \section{Aufgaben}
    Diese Aufgaben sollen Ihnen dabei helfen, ein tieferes Verständnis für Linux und seine Möglichkeiten zu entwickeln.

    \subsection{Aufgabe 1}
    Was ist der Linux-Kernel und welche Aufgaben erfüllt er?

    \subsection{Aufgabe 2}
    Informieren Sie sich über Container-Technologien wie Docker.
    Welche Vorteile haben diese gegenüber virtuellen Maschinen? \\ \\
    Erläutern Sie ebenfalls den Zusammenhang zwischen Linux-Kernel und Container.
    Starten Sie einen einfachen Container und untersuchen Sie dessen Ressourcenverbrauch.

    \subsection{Aufgabe 3}
    Vergleichen Sie verschiedene Dateisysteme wie ext4, Btrfs oder ZFS.

    \subsection{Aufgabe 4}
    Recherchieren Sie Sicherheitsmechanismen wie SELinux oder AppArmor.
    Welche Probleme sollen diese lösen?

    \subsection{Aufgabe 5}
    (Optional) Finden Sie heraus, wie man einen eigenen Kernel konfiguriert oder kompiliert.

\end{document}